% Ana - Ejercicio 2.b: SIG SEMARNAT.
% TODO: desarrollar después de acordar el reparto.

\section{Ejercicio 2.b: Traducción del modelo E/R al modelo relacional}

Reglas de Traducción Aplicadas:
\begin{itemize}
    \item Entidad fuerte: se convierte en tabla con sus atributos.
    \item Relaciones M:N: se convierten en tabla con las llaves de las entidades relacionadas más los atributos de la relación.
    \item Relaciones 1:N: la entidad del lado N incluye la llave de la entidad del lado 1 más los atributos de la relación.
    \item Especialización-generalización: se crea una tabla para la superentidad y una tabla para cada subentidad con los atributos propios más la llave de la superentidad.
    \item Clasificación de entidades: se crea una tabla para la superentidad y una tabla para cada subentidad con los atributos propios más la llave de la superentidad.
    \item Relación recursiva: se crea una tabla para la relación. Se crearon dos llaves distintas que puntan a la misma entidad para diferenciar los roles de la relación.
\end{itemize}

Decisiones de Diseño
\begin{enumerate}
    \item \textbf{Relación Alimentar}, la cual es recursiva ya que un Río puede alimentar a otro río. Se creó una tabla para la relación con dos llaves foráneas que apuntan a la misma entidad Río, diferenciando los roles de la relación.
\end{enumerate}
